\section{Visión general completa de los modelos de difusión}

\subsection{Simetría entre el proceso forward y el inverso}

El marco completo de los modelos de difusión puede resumirse mediante la siguiente estructura simétrica:

\begin{table}[H]
\centering
\begin{tabular}{lll}
\toprule
 & \textbf{Proceso forward (adición de ruido)} & \textbf{Proceso inverso (eliminación de ruido)} \\
\midrule
Dirección & $x_0 \to x_1 \to \cdots \to x_T$ & $x_T \to x_{T-1} \to \cdots \to x_0$ \\
Distribución & $q(x_t | x_{t-1})$, predefinida & $p_\theta(x_{t-1} | x_t)$, por aprender \\
¿Entrenamiento? & No & Sí \\
Punto final & $x_T \approx \mathcal{N}(0, I)$ & $x_0 \sim p_{\text{data}}$ \\
Función & Generar muestras ruidosas durante el entrenamiento & Generar muestras de datos durante la inferencia \\
\bottomrule
\end{tabular}
\caption{Comparación entre el proceso forward y el inverso}
\end{table}

\subsection{¿Por qué son efectivos los modelos de difusión?}

\begin{knowledgebox}{Comprensión de los modelos de difusión desde la perspectiva de VAE}
Los modelos de difusión pueden entenderse como un tipo especial de VAE jerárquico, con las siguientes características clave:
\begin{itemize}
    \item El codificador (proceso forward) es fijo y no requiere aprendizaje.
    \item Todas las capas tienen la misma dimensión.
    \item La transición en cada paso es una perturbación gaussiana simple.
    \item Mediante suficientes pasos, el producto de transiciones gaussianas simples puede expresar distribuciones muy complejas.
\end{itemize}
Este diseño logra un equilibrio elegante entre simplicidad y capacidad expresiva.
\end{knowledgebox}

\subsection{Próximo contenido}

Esta clase establece el marco conceptual de DDPM; la siguiente clase (Lecture 03: DDPM Parte 2) profundizará en ello:

\begin{itemize}
    \item Derivación detallada del término de consistencia $L_{t-1}$
    \item Expresión gaussiana cerrada para $q(x_{t-1}|x_t, x_0)$
    \item Simplificación desde la predicción de la media hasta la \textbf{predicción de ruido} $\epsilon_\theta$
    \item Objetivo final de entrenamiento de DDPM: una pérdida MSE ponderada simple
    \item Algoritmos reales de entrenamiento y muestreo
\end{itemize}

\subsection{Resumen de la sección}

Los modelos de difusión logran una potente capacidad generativa dentro del marco de VAE jerárquicos mediante un proceso forward fijo (adición de ruido) y un proceso inverso aprendido (eliminación de ruido). El diseño ingenioso del proceso forward hace que varios términos del ELBO se cumplan automáticamente o puedan calcularse eficientemente, enfocando el entrenamiento en el término de consistencia central.
